% TOP_PROBLEM: {"id": 30007482, "title": "One rational matrix hitting point for noncommutative formulas", "category_id": 15, "category": {"id": 15, "name": "computer_science", "display_name": "Computer Science", "description": "Computational complexity, algorithms, and theoretical CS.", "slug": "computer-science", "order_index": 15, "created_at": "2026-07-31T15:26:25.671Z"}, "status": "open"}
% FIRST_POSTED: 2026-10-06
% ATTEMPT_STATUS: solved
% OPENAI_CLAIM: {"schema_version":1,"source_repo":"https://github.com/openai/math","source_commit":"adc7f1241b42e322a6451854ab7e4b4c146bf78a","release_date":"2026-10-06","match":"full","resolution":"proved","independently_reviewed":false,"second_pass":"agreed","families":[{"family":"116","title":"Uniform black-box noncommutative identity testing across characteristics","subject":"Theoretical computer science","manuscripts":[{"title":"One Rational Matrix Hitting Point for Noncommutative Formulas","pdf_path":"preprints/One-Rational-Matrix-Hitting-Point-for-Noncommutative-Formulas-September-24-2026/One-Rational-Matrix-Hitting-Point-for-Noncommutative-Formulas-September-24-2026.pdf","readme_path":"preprints/One-Rational-Matrix-Hitting-Point-for-Noncommutative-Formulas-September-24-2026/README.md","date":"2026-09-24","theorem_ref":"thm:main (introduction.tex:24)","statement_sources":[{"path":"preprints/One-Rational-Matrix-Hitting-Point-for-Noncommutative-Formulas-September-24-2026/build/introduction.tex","line":24,"label":"thm:main"}]}],"lean":{"doc_path":"lean/docs/116.md","comparators":["lean/ComparatorChallenges/FormulaHitting.json","lean/ComparatorChallenges/FormulaHitting.lean","lean/ComparatorChallenges/RationalHitting.json","lean/ComparatorChallenges/RationalHitting.lean"],"declarations":["OAI.RationalHitting.main"],"covers_main_theorem":false},"reasoning_trace":null}]}
\documentclass[11pt]{article}
\usepackage[margin=25mm]{geometry}
\usepackage{fontspec}
\setmainfont{Latin Modern Roman}
\usepackage{amsmath,amssymb,mathtools}
\usepackage{unicode-math}
\setmathfont{Latin Modern Math}
\usepackage{enumitem}
\usepackage{xurl,hyperref,bookmark}
\hypersetup{colorlinks=true,urlcolor=blue}
\setcounter{secnumdepth}{0}
\setlength{\parindent}{0pt}
\setlength{\parskip}{5pt}
\setlength{\emergencystretch}{3em}
\begin{document}
\section*{One rational matrix hitting point for noncommutative formulas}

\subsection{OpenAI's claim}
OpenAI presents the following claims relevant to this question, with manuscript numbers following paper order in Sources. In manuscript 1 (family 116; 2026-09-24), the asserted result is that one rational matrix tuple of dimension $O(ns^2)$ detects every nonzero size-$s$ division-free noncommutative formula and is constructible in polynomial bit time.

\subsection{Mathematical statement of the claimed result}
\paragraph{Manuscript 1.}
For integers $n,s\ge1$, the claim gives a deterministic algorithm constructing one rational matrix tuple $\mathbf T_{n,s}\in M_d(\mathbb Q)^n$ in polynomial bit time, with $d\le2ns^2$, such that
\[
 \forall F\text{ of characteristic }0\ \forall0\ne f\in\mathsf{Form}_{n,s}(F):\quad f(\mathbf T_{n,s})\ne0.
\]
Here $\mathsf{Form}_{n,s}$ denotes noncommutative arithmetic formulas in $n$ variables of size at most $s$, and matrices are viewed in $F$ via $\mathbb Q\hookrightarrow F$. Entry numerators and denominators have bit length $O(\log(ns+1)+ns^2\log(n+1))$. The stated dimensions are $2ns^2$ if $ns\ge2$ and one if $n=s=1$.

\subsection{Outline of the argument}
For manuscript 1 (family 116; 2026-09-24), iterated integrals encode free words faithfully inside differential systems. A vanishing-order bound limits how many Taylor coefficients are needed; those coefficients form the finite matrix space used for the universal substitution.

\subsection{Scope relative to this problem}
The result supplies the conclusion in this problem with its displayed hypotheses and quantifiers. The broader statements and applications in the manuscript are kept separate from the particular question above. The supporting result is manuscript 1 (family 116; 2026-09-24).

\subsection{Formal verification}
For manuscript 1 (family 116; 2026-09-24), the documented Lean scope is the universal matrix hitting property over characteristic-zero fields; the construction-time and $O(ns^2)$ dimension estimates are not separately encoded in the selected theorem.

This coverage does not establish a formal proof of the complete problem as stated here.

The linked family Comparator checks are \texttt{FormulaHitting.lean}, \texttt{RationalHitting.lean}.

Their named declarations include \texttt{OAI.NCHitting.universal\_hitting}, \texttt{OAI.RationalHitting.main}.

This site has not executed the Lean code. Described statement coverage is not an independent proof check.

\subsection{Status caveat}
The manuscripts were produced by an internal OpenAI model. OpenAI cautions that results without formalization can contain mistakes. This site has not independently checked the proofs; the label records the scope of OpenAI's claim.

\subsection{Sources}
\begin{itemize}
\item OpenAI, \emph{One Rational Matrix Hitting Point for Noncommutative Formulas}, 2026-09-24; citation key \texttt{OAI:One-Rational-Matrix-Hitting-Point-for-Noncommutative-Formulas-September-24-2026}. \href{https://github.com/openai/math/blob/adc7f1241b42e322a6451854ab7e4b4c146bf78a/preprints/One-Rational-Matrix-Hitting-Point-for-Noncommutative-Formulas-September-24-2026/One-Rational-Matrix-Hitting-Point-for-Noncommutative-Formulas-September-24-2026.pdf}{PDF (GitHub)}; \href{https://github.com/openai/math/raw/adc7f1241b42e322a6451854ab7e4b4c146bf78a/preprints/One-Rational-Matrix-Hitting-Point-for-Noncommutative-Formulas-September-24-2026/One-Rational-Matrix-Hitting-Point-for-Noncommutative-Formulas-September-24-2026.pdf}{PDF (direct)}; \href{https://github.com/openai/math/blob/adc7f1241b42e322a6451854ab7e4b4c146bf78a/preprints/One-Rational-Matrix-Hitting-Point-for-Noncommutative-Formulas-September-24-2026/README.md}{Manuscript page and citation}; \href{https://github.com/openai/math/blob/main/preprints/One-Rational-Matrix-Hitting-Point-for-Noncommutative-Formulas-September-24-2026/One-Rational-Matrix-Hitting-Point-for-Noncommutative-Formulas-September-24-2026.pdf}{Latest version}; Statement source: \href{https://github.com/openai/math/blob/adc7f1241b42e322a6451854ab7e4b4c146bf78a/preprints/One-Rational-Matrix-Hitting-Point-for-Noncommutative-Formulas-September-24-2026/build/introduction.tex\#L24}{thm:main}.
\item Family 116: \href{https://github.com/openai/math/blob/adc7f1241b42e322a6451854ab7e4b4c146bf78a/lean/docs/116.md}{Lean scope}.
\item \href{https://github.com/openai/math/blob/adc7f1241b42e322a6451854ab7e4b4c146bf78a/lean/ComparatorChallenges/FormulaHitting.json}{Comparator FormulaHitting.json}.
\item \href{https://github.com/openai/math/blob/adc7f1241b42e322a6451854ab7e4b4c146bf78a/lean/ComparatorChallenges/FormulaHitting.lean}{Comparator FormulaHitting.lean}.
\item \href{https://github.com/openai/math/blob/adc7f1241b42e322a6451854ab7e4b4c146bf78a/lean/ComparatorChallenges/RationalHitting.json}{Comparator RationalHitting.json}.
\item \href{https://github.com/openai/math/blob/adc7f1241b42e322a6451854ab7e4b4c146bf78a/lean/ComparatorChallenges/RationalHitting.lean}{Comparator RationalHitting.lean}.
\item \href{https://github.com/openai/math/blob/adc7f1241b42e322a6451854ab7e4b4c146bf78a/CONTENTS.md}{OpenAI catalogue entry}.
\end{itemize}
\end{document}
